\section{The no-go theorems and where associativity enters}
\label{sec:nogo}

\subsection{What the theorems say}

A \emph{generator} of a continuous symmetry is the infinitesimal transformation: the derivative, at
the identity, of a one-parameter family of transformations. The generators of a group of
transformations form a vector space with a \emph{bracket}, the generator of the commutator of two
infinitesimal transformations; for a Lie group this is a \emph{Lie algebra}: the bracket is bilinear,
antisymmetric, and obeys the Jacobi identity, i.e.\ its \emph{Jacobiator}
\begin{equation}
\mathrm{Jac}_{\rm LL}(\mathbf x,\mathbf y,\mathbf z)=[[\mathbf x,\mathbf y],\mathbf z]+[[\mathbf y,\mathbf z],\mathbf x]+[[\mathbf z,\mathbf x],\mathbf y]
\label{eq:jacobiator}
\end{equation}
vanishes for all generators $\mathbf x$, $\mathbf y$, $\mathbf z$. (In this section $\mathbf x$,
$\mathbf y$, $\mathbf z$ stand for generators; from Sec.~\ref{sec:bioct} on they are bi-octonions.)
The Lorentz algebra $\mathfrak{so}(3,1)$ has six generators, three rotations and three boosts; adding
the four translations gives the Poincar\'e algebra.

The no-go theorems state, under hypotheses listed below, that a symmetry algebra containing the
Poincar\'e algebra and an internal algebra is their direct sum:
\begin{itemize}
\item McGlinn~\cite{McGlinn1964}: in a Lie algebra that contains the Poincar\'e algebra and an
internal algebra, and in which the translations act nilpotently, the internal generators commute with
the Poincar\'e generators.
\item O'Raifeartaigh~\cite{ORaifeartaigh1965}: if a finite-dimensional Lie algebra contains the
Poincar\'e algebra and acts by a unitary representation in which the squared mass has an isolated
eigenvalue, the members of a multiplet all have the same mass; the algebra cannot split masses.
\item Coleman and Mandula~\cite{ColemanMandula1967} (in any number of space dimensions with one time,
Pelc and Horwitz~\cite{PelcHorwitz1997}): the most general Lie group of symmetries of a non-trivial
scattering matrix of a relativistic theory with finitely many particles below any mass, analytic
amplitudes and a mass gap is, locally, the direct product of the Poincar\'e group and a compact
internal group.
\item Haag, {\L}opusza\'nski and Sohnius~\cite{HLS1975} (see also~\cite{WeinbergQFT3}): if the
generators are allowed to form a \emph{Lie superalgebra}, in which some generators (called odd) have
an anticommutator instead of a commutator, the odd generators are spinors whose anticommutators are
translations (and central charges): supersymmetry.
\end{itemize}

\subsection{The hypotheses}

We label the hypotheses that the proofs use:
\begin{itemize}
\item[H-grp.] The symmetry transformations form a group: their composition is associative, there is
an identity, and every transformation has an inverse.
\item[H-op.] They act on the states by linear (unitary or antiunitary) operators, so that the
composition of transformations is the product of operators; particles are irreducible unitary
representations of the Poincar\'e group~\cite{Wigner1939}.
\item[H-Lie.] The group is a Lie group; its generators form a Lie algebra.
\item[H-gr.] (Haag, {\L}opusza\'nski and Sohnius only) The generators form a Lie superalgebra, whose
graded Jacobi identity is the Jacobi identity with the signs of the anticommutators.
\item[H-dyn.] The dynamical hypotheses on the scattering matrix: it is non-trivial, the amplitudes are
analytic, there are finitely many particles below any mass, and there is a mass gap.
\end{itemize}
Every theorem above uses H-Lie (or H-gr), and the theorems on the scattering matrix use H-op and
H-dyn as well.

\subsection{Where associativity enters}

H-grp states associativity. H-op implies it, because the product of operators is associative. And
H-Lie follows from either. Take any three elements $\mathsf O_1$, $\mathsf O_2$, $\mathsf O_3$ of an
algebra (for example operators), their commutator $[\mathsf O_1,\mathsf O_2]=\mathsf O_1\mathsf
O_2-\mathsf O_2\mathsf O_1$, and their \emph{associator}
$[\mathsf O_1,\mathsf O_2,\mathsf O_3]=(\mathsf O_1\mathsf O_2)\mathsf O_3-\mathsf O_1(\mathsf O_2\mathsf O_3)$.
Expanding the twelve products of the Jacobiator and sorting them by the order of the factors gives
\begin{align}
&\mathrm{Jac}_{\rm LL}(\mathsf O_1,\mathsf O_2,\mathsf O_3)\nonumber\\
&\quad=[\mathsf O_1,\mathsf O_2,\mathsf O_3]+[\mathsf O_2,\mathsf O_3,\mathsf O_1]+[\mathsf O_3,\mathsf O_1,\mathsf O_2]\nonumber\\
&\quad\phantom{=}-[\mathsf O_2,\mathsf O_1,\mathsf O_3]-[\mathsf O_1,\mathsf O_3,\mathsf O_2]-[\mathsf O_3,\mathsf O_2,\mathsf O_1]:
\label{eq:jacassoc}
\end{align}
the associators of the even orderings of the three elements minus those of the odd orderings. For an
associative product (operators) every associator vanishes, so the Jacobi identity holds identically (check A-6 verifies Eq.~(\ref{eq:jacassoc}) for random integer
matrices). Equation~(\ref{eq:jacassoc}) holds for any product, associative or not, and it is the hinge
of this paper: in a non-associative algebra the commutator need not obey the Jacobi identity, and the
generators of its symmetries need not form a Lie algebra.

\subsection{What a counterexample can be}

A counterexample in which all hypotheses of a theorem hold does not exist; a search for one would end
in an error. We therefore construct structures in which H-grp, H-op and H-Lie (and H-gr) fail
\emph{because associativity fails}, while Lorentz covariance is kept exactly, and in which the
conclusion of the theorems---no non-trivial union of Lorentz and internal symmetry---fails. The
dynamical hypotheses H-dyn concern the scattering matrix of a quantum theory; because our structures
have no realization by operators (Theorem~R, Sec.~\ref{sec:bioct}), they cannot act on the states of an
ordinary quantum theory, and H-dyn is not tested. A quantum mechanics over non-associative algebras has
been studied since Jordan, von Neumann and Wigner~\cite{JordanNeumannWigner1934} and, for the octonions,
by G\"unaydin, Piron and Ruegg~\cite{GunaydinPironRuegg1978}; we do not use it.
